\section{Sharp interpolation and sparse nonabelian tests}\label{sec:sparse-sampling}

\paragraph{Status and attribution.}
Fourth-power sampling is already present in Sandi Setiawan's
\href{https://www.preprints.org/manuscript/202611.0738}{November 2026 preprint},
whose abstract explicitly states the fourth-power square-root scale.
We make no novelty claim for that scale. This section gives a self-contained
version for arbitrary positive ramp measures, an exact interpolation loss,
a sharpness example within that class, and a centering-free specialization
to nonabelian Dedekind zeta functions. The latter is a reduction of an open
estimate, not a proof of that estimate.

Let $\mu$ be a positive locally finite Borel measure on $[1,\infty)$,
let $A\ge0$, and let $c\in\mathbb R$. Put
\[
 S(x)=\int_{[1,x]}\left(\frac1u-\frac1x\right)d\mu(u),
 \qquad R(x)=S(x)-A\log x-c\quad(x\ge1).
\]
No convergence assumption on $R$ is needed for the interpolation argument.
For prime weights, $\mu=\sum_p w(p)\log p\,\delta_p$ with $w(p)\ge0$.
For the source repository's remainder, use the nonnegative prime-power
weights $\lambda_{d,q}(n)$ and $A=1/\varphi(q)$.

\begin{lemma}[Exact interpolation loss]\label{lem:sampling-jensen}
Suppose $1\le x_0<x_1$ and $x\in[x_0,x_1]$. Define
\[
 t_i=x_i^{-1},\qquad t=x^{-1},\qquad
 \lambda=\frac{t-t_1}{t_0-t_1}\in[0,1].
\]
Then
\begin{align}
 R(x)&\le\lambda R(x_0)+(1-\lambda)R(x_1)+A J_\lambda(t_0,t_1),
 \label{eq:sampling-interpolation}\\
 J_\lambda(t_0,t_1)
 &=\log\bigl(\lambda t_0+(1-\lambda)t_1\bigr)
       -\lambda\log t_0-(1-\lambda)\log t_1\nonumber\\
 &\le\frac18\log^2(x_1/x_0).
 \label{eq:sampling-loggap}
\end{align}
The exact maximum of $J_\lambda$ over $\lambda\in[0,1]$ is
\[
 J(r)=\log\frac{r-1}{\log r}-1+\frac{\log r}{r-1},
 \qquad r=x_1/x_0>1.
\]
As $r\downarrow1$, $J(r)\sim\frac18\log^2r$.
\end{lemma}

\begin{proof}
The function
\[
 H(t)=S(1/t)=\int_{[1,\infty)}(u^{-1}-t)_+\,d\mu(u),\qquad t>0,
\]
is convex, since each integrand is the maximum of two affine functions
and the measure is positive. Only $u\le1/t$ contributes, so every integral
is finite. Apply convexity to $H(t)$ and substitute
$H(t)=R(1/t)-A\log t+c$. This gives
\eqref{eq:sampling-interpolation}.

For the bound, set $y=\log(t_0/t_1)$ and
$f_\lambda(y)=\log(1-\lambda+\lambda e^y)-\lambda y$.
The cases $\lambda=0,1$ are immediate. Otherwise
\[
 f_\lambda(0)=f'_\lambda(0)=0,
 \qquad f''_\lambda(y)=p(y)(1-p(y))\le\tfrac14,
 \quad p(y)=\frac{\lambda e^y}{1-\lambda+\lambda e^y}.
\]
Two integrations give $f_\lambda(y)\le y^2/8$.
For the exact maximum, write $u=t/t_1\in[1,r]$. The gap is
$\log u-(u-1)\log r/(r-1)$, maximized at
$u=(r-1)/\log r$. Substitution proves the formula. Taylor expansion
at $r=1$ gives the asymptotic; it also follows by the preceding
second-derivative calculation at $\lambda=1/2$.
\end{proof}

\begin{theorem}[Nonuniform sampling]\label{thm:sampling-general}
Let $\theta>0$ and $1\le x_1<x_2<\cdots\to\infty$. Suppose
\[
 \log(x_{j+1}/x_j)=O(x_j^{-\theta/2}).
\]
Then the following are equivalent:
\[
 \exists C\ge0\ \forall x\text{ sufficiently large},\quad
 R(x)\le Cx^{-\theta};
\]
\[
 \exists C_0\ge0\ \forall j\text{ sufficiently large},\quad
 R(x_j)\le C_0x_j^{-\theta}.
\]
More quantitatively, if the sampled inequality holds and
$\log^2(x_{j+1}/x_j)\le Kx_j^{-\theta}$, then for
$x\in[x_j,x_{j+1}]$,
\[
 R(x)\le(C_0+AK/8)x_j^{-\theta}
 \le(C_0+AK/8)(x_{j+1}/x_j)^\theta x^{-\theta}.
\]
When $A>0$, the scale of the required gaps is
$x_{j+1}-x_j=O(x_j^{1-\theta/2})$.
\end{theorem}

\begin{proof}
The continuous inequality implies its restriction to the samples.
For the reverse implication, both endpoint values in
\eqref{eq:sampling-interpolation} are at most $C_0x_j^{-\theta}$;
use \eqref{eq:sampling-loggap}. The ratio $x_{j+1}/x_j$ tends to one,
so it is bounded. The assertion about gaps follows from
$e^v-1\sim v$ as $v\to0$.
\end{proof}

\begin{corollary}[Fourth powers]\label{cor:sampling-quartic}
For every positive ramp measure as above,
\[
 R(x)=O_+(x^{-1/2})
 \quad\Longleftrightarrow\quad
 R(n^4)=O_+(n^{-2}),
\]
where $O_+$ means an eventual upper bound, with no lower bound asserted.
If $R(n^4)\le Cn^{-2}$ for $n\ge n_0$ and $C\ge0$, then
\[
 R(x)\le(C+8A)x^{-1/2}\qquad(x\ge n_0^4).
\]
For the polynomial grid $x_j=j^r$, the general theorem applies whenever
$r\theta\le2$.
\end{corollary}

\begin{proof}
Here $\log(x_{n+1}/x_n)=4\log(1+1/n)\le4/n$, so the interpolation
loss is at most $2A/n^2$. Because $C\ge0$ and the square root is
concave, the endpoint contribution is
$C[\lambda\sqrt{t_0}+(1-\lambda)\sqrt{t_1}]\le C\sqrt t$.
For $x\in[n^4,(n+1)^4]$ one has $n^{-2}\le4x^{-1/2}$.
Thus the total loss is at most $8Ax^{-1/2}$. In general the logarithmic gap is
$r\log(1+1/j)\sim r/j$, while $x_j^{-\theta}=j^{-r\theta}$.
\end{proof}

There is a sharper asymptotic statement. If
$L=\limsup_{n\to\infty}n^2R(n^4)$ is finite, then
\[
 L\le\limsup_{x\to\infty}\sqrt x\,R(x)\le L+2A.
\]
Indeed, for $x\in[n^4,(n+1)^4]$ the normalized endpoint weights
$\sqrt x/n^2$ and $\sqrt x/(n+1)^2$ tend uniformly to one.
Multiplying \eqref{eq:sampling-interpolation} by $\sqrt x$ bounds the
endpoint term by $L+o(1)$, with an arbitrarily small fixed slack,
and bounds the Jensen loss by
$2A(n+1)^2\log^2(1+1/n)=2A+o(1)$. The first inequality is restriction
to the sampled subsequence. The same argument handles $L=-\infty$
by first imposing an arbitrary negative bound on the sampled terms.

\begin{proposition}[Sharpness of the positivity argument]\label{prop:sampling-sharp}
Fix $A>0$, $r>0$, and nodes $x_j=j^r$. There is a positive measure
$\mu$ supported on these nodes for which $c=0$,
\[
 R(x_j)=0\quad\text{for every }j,\qquad R(x)\longrightarrow0,
\]
but at suitable $x_j^*\in(x_j,x_{j+1})$,
\[
 R(x_j^*)=A J((1+1/j)^r)\sim\frac{Ar^2}{8j^2}.
\]
Consequently $x^{\theta}R(x)$ is unbounded above if $r\theta>2$.
If $r$ is an integer, the measure is supported on positive integers.
At the critical case $r=4$, $\theta=1/2$, the asymptotic loss $2A$
cannot be improved using positivity alone.
\end{proposition}

\begin{proof}
Write $t_j=j^{-r}$ and define $H(t)$ to be the piecewise affine
interpolant of $-A\log t$ on the nodes $t_j$, with $H(t)=0$ for
$t\ge1$. It is convex. On $[t_{j+1},t_j]$ its slope is $-D_j$, where
\[
 D_j=\frac{A\log(x_{j+1}/x_j)}{x_j^{-1}-x_{j+1}^{-1}}.
\]
The $D_j$ strictly increase. Set $D_0=0$ and
\[
 \mu=\sum_{j\ge1}(D_j-D_{j-1})\delta_{x_j}.
\]
This is positive and locally finite. The function
$\sum_j(D_j-D_{j-1})(x_j^{-1}-t)_+$ has the same piecewise slopes
and the same value at $t=1$ as $H$, so it equals $H$.
Therefore $S(x_j)=A\log x_j$ and $R(x_j)=0$.
Between nodes the remainder is exactly the Jensen gap in
Lemma~\ref{lem:sampling-jensen}. Its maximum occurs at
\[
 x_j^*=\frac{\log(x_{j+1}/x_j)}{x_j^{-1}-x_{j+1}^{-1}}
       \sim j^r.
\]
The asserted asymptotic follows from that lemma. It also implies
$R(x)\to0$, since the maximum on the $j$th interval is $O(j^{-2})$.
Multiplying by $(x_j^*)^{\theta}$ proves the divergence assertion.
For $r=4$ the limiting normalized maximum is $2A$.
\end{proof}

This counterexample is deliberately a statement about what positivity
and interpolation can prove. Its masses satisfy
$\mu(\{j^r\})\sim Arj^{r-1}$; they are not von Mangoldt coefficients.
It does not establish that a sparser arithmetic sampling criterion is
false. Analytic continuation or additional arithmetic structure could
give stronger sampling results.

\subsection{Eliminating the unknown centering constant}

Suppose now that $R(x)\to0$ and let $q\ge2$ be an integer. For
the quartic grid define the explicitly finite statistic
\begin{align}
 B_q(n)&=4A\log q-\bigl(S((qn)^4)-S(n^4)\bigr)\label{eq:sampling-B}\\
 &=R(n^4)-R((qn)^4).\nonumber
\end{align}
The centering constant $c$ cancels. With $X=n^4$ and $Q=q^4$,
\begin{align}
 S(QX)-S(X)
 &=\frac{Q-1}{QX}\,\mu([1,X])
   +\int_{(X,QX]}\left(\frac1u-\frac1{QX}\right)d\mu(u).
 \label{eq:sampling-finite-increment}
\end{align}
Thus only the weights at $u\le(qn)^4$ are used.

\begin{proposition}[One-sided finite certificate]\label{prop:sampling-certificate}
If $B_q(n)\le Cn^{-2}$ for all sufficiently large integers $n$,
with $C\ge0$, then
\[
 R(n^4)\le\frac{C}{1-q^{-2}}n^{-2}
\]
for those integers, and hence $R(x)=O_+(x^{-1/2})$.
Conversely, the two-sided estimate $R(x)=O(x^{-1/2})$ implies
$B_q(n)=O(n^{-2})$.
\end{proposition}

\begin{proof}
For every integer $N\ge0$,
\[
 R(n^4)-R((q^{N+1}n)^4)=\sum_{k=0}^{N}B_q(q^kn)
 \le\frac{C}{n^2}\sum_{k=0}^{N}q^{-2k}.
\]
Let $N\to\infty$ and use $R(x)\to0$. Then apply
Corollary~\ref{cor:sampling-quartic}. The converse follows directly
from \eqref{eq:sampling-B} and a two-sided bound at both endpoints.
\end{proof}

The converse here deliberately assumes a \emph{two-sided} estimate.
Subtracting two upper bounds does not bound their difference.
For zeta-function applications RH supplies the two-sided estimate, so
this distinction does not weaken the resulting equivalence.

\subsection{A nonabelian criterion using polynomial root counts}

Let $f\in\mathbb Z[T]$ be monic and irreducible, let
$K=\mathbb Q[T]/(f)$, and omit primes dividing $\operatorname{disc}(f)$.
For every other prime $p$, let
\[
 r_f(p)=\#\{a\in\mathbb F_p:f(a)=0\}.
\]
It counts the degree-one prime ideals of $K$ above $p$ and lies between
zero and $[K:\mathbb Q]$. Define
\[
 S_f(x)=\sum_{\substack{p\le x\\p\nmid\operatorname{disc}(f)}}
       r_f(p)\log p\left(\frac1p-\frac1x\right),
 \qquad R_f(x)=S_f(x)-\log x-c_f,
\]
where $c_f$ is the unique constant making $R_f(x)\to0$.
Its existence follows from the fixed-field prime ideal theorem:
if $\theta_f(x)=\sum_{p\le x}r_f(p)\log p$, then
\[
 \theta_f(x)=x+O_f(xe^{-b_f\sqrt{\log x}}),
 \qquad
 c_f=\int_1^\infty(\theta_f(t)-t)t^{-2}\,dt.
\]
In this estimate a possible exceptional zero of the fixed field is
absorbed for sufficiently large $x$. Higher residue degrees and ideal
powers cost $O_f(\sqrt x\log^2(2x))$, which is also absorbed.

\begin{theorem}[A finite quartic criterion for any number field]
\label{thm:sampling-numberfield}
With the notation above, set
\begin{align}
 B_f(n)=\log16
 &-\frac{15}{16n^4}
       \sum_{\substack{p\le n^4\\p\nmid\operatorname{disc}(f)}}r_f(p)\log p
 \nonumber\\[-2pt]
 &-\sum_{\substack{n^4<p\le16n^4\\p\nmid\operatorname{disc}(f)}}
       r_f(p)\log p\left(\frac1p-\frac1{16n^4}\right).
 \label{eq:sampling-polynomial-statistic}
\end{align}
Then
\[
 \operatorname{RH}(\zeta_K)
 \quad\Longleftrightarrow\quad
 \sup_{n\ge1}n^2B_f(n)<\infty.
\]
Equivalently, RH holds if and only if
$\sup_{n\ge1}n^2R_f(n^4)<\infty$.
The bound on $B_f$ is one-sided; no unspecified centering constant
appears in its definition.
\end{theorem}

\begin{proof}
We first justify the continuous analytic criterion. For $\Re s>1$,
\[
 D_f(s)=\sum_{p\nmid\operatorname{disc}(f)}r_f(p)(\log p)p^{-s}
       =-\frac{\zeta'_K(s)}{\zeta_K(s)}+H_f(s),
\]
where $H_f$ is holomorphic in $\Re s>1/2$. Indeed the omitted
prime-power coefficients have absolute value at most
$[K:\mathbb Q]\log p$, and the series with exponent $p^{-js}$,
$j\ge2$, converges locally uniformly there. The finitely many
ramified local factors are holomorphic for $\Re s>0$.
Termwise integration gives
\[
 T_f(s)=\int_1^\infty R_f(x)x^{-s}\,dx
       =\frac{D_f(s)}{s(s-1)}-\frac1{(s-1)^2}-\frac{c_f}{s-1}.
\]
The apparent singularity at one is removable by the definition of
$c_f$. Every zero $\rho$ of $\zeta_K$ with $\Re\rho>1/2$ supplies
a simple pole with residue
\[
 \frac{m_\rho}{\rho(1-\rho)},
\]
where $m_\rho>0$ is its multiplicity. In particular each real pole
in $(1/2,1)$ has positive residue, and there are no poles on the real
line at or above one.

If $R_f(x)\le Cx^{-1/2}$ eventually, fix $\beta>1/2$. The function
$Cx^{\beta-1}-R_f(x)$ is eventually nonnegative after increasing $C$.
Its tail Mellin transform equals $C/(s-\beta)-T_f(s)$ plus an entire
function. A convergence abscissa to the right of $\beta$ would, by
Landau's lemma for a nonnegative integrand, force a real singularity
there. Such a singularity is either removable or a simple pole with
negative residue, the latter impossible for an integral tending to
a nonnegative value or $+\infty$ from the right. Thus its transform
is holomorphic in $\Re s>\beta$ and $T_f$ has no poles there.
As $\beta\downarrow1/2$, there are no zeros to the right of the
critical line. The Dedekind functional equation supplies the other
half of the critical strip.

In the forward direction, under RH the integrated explicit formula
gives
\[
 \int_0^x(\psi_K(t)-t)\,dt=O_K(x^{3/2}).
\]
This form is given, with explicit constants, by
\href{https://arxiv.org/abs/1312.4465}{Greni\'e and Molteni,
\emph{Explicit smoothed prime ideals theorems under GRH}}.
The difference $\psi_K(t)-\theta_f(t)=O_f(\sqrt t\log^2(2t))$
can be sharpened here to $O_f(\sqrt t)$: for $j=2$ use
$\sum_{p\le\sqrt t}\log p=O(\sqrt t)$; the sum over $j\ge3$
is $O_f(t^{1/3}\log(2t))$, and finitely many bad primes contribute
$O_f(\log(2t))$. Therefore
$A_f(x)=\int_1^x(\theta_f(t)-t)\,dt=O_f(x^{3/2})$.
Integration by parts in
$R_f(x)=-\int_x^\infty(\theta_f(t)-t)t^{-2}\,dt$ gives
\[
 R_f(x)=\frac{A_f(x)}{x^2}-2\int_x^\infty\frac{A_f(t)}{t^3}\,dt
       =O_f(x^{-1/2}).
\]
This proves the continuous equivalence, including the exact exponent.

Apply Corollary~\ref{cor:sampling-quartic} to the positive measure
$\sum_p r_f(p)\log p\,\delta_p$, with $A=1$.
Finally \eqref{eq:sampling-polynomial-statistic} is $B_2(n)$ of
\eqref{eq:sampling-B}. Proposition~\ref{prop:sampling-certificate}
proves the reverse implication from its upper bound, and the
two-sided estimate just proved gives the forward implication.
Finitely many initial terms are finite, so eventual boundedness and
the displayed supremum condition are equivalent.
\end{proof}

\paragraph{The first nonabelian example.}
Take $f(T)=T^3-T-1$, let $\alpha$ be a root, and set
$K=\mathbb Q(\alpha)$. The rational-root test proves irreducibility,
and $\operatorname{disc}(f)=4-27=-23$; hence the Galois closure $M$
has group $S_3$, with quadratic subfield $\mathbb Q(\sqrt{-23})$.
Only $p=23$ is omitted from the prime sum. For every other prime,
\[
 r_f(p)=\deg\gcd(T^p-T,T^3-T-1)\quad\text{in }\mathbb F_p[T],
\]
so it is computable by modular polynomial arithmetic. The possible
factorization types are three linear factors, one linear factor times
an irreducible quadratic, and an irreducible cubic; their root counts
are respectively $3,1,0$. The corresponding Frobenius characters are
\[
 \begin{array}{c|ccc}
 \text{type}&111&12&3\\ \hline
 r_f&3&1&0\\
 \operatorname{sgn}&1&-1&1\\
 \operatorname{tr}\rho&2&0&-1
 \end{array}
\]
with $r_f=1+\operatorname{tr}\rho$ for the standard two-dimensional
representation $\rho$ of $S_3$. The uniform upper bound in
Theorem~\ref{thm:sampling-numberfield} would therefore prove RH for
this nonabelian cubic field. The repository's Dirichlet GRH does not
supply that upper bound: the root-count character contains $\rho$.

For the Galois closure, the analogous statistic uses the indicator
$r_f(p)=3$ and replaces $\log16$ by $(\log16)/6$. Its boundedness is
equivalent to $\operatorname{RH}(\zeta_M)$ by the completely split
criterion in the prior reports. Under Dirichlet GRH these two target
RH statements are equivalent, since
\[
 \zeta_K=\zeta L(s,\rho),\qquad
 \zeta_M=\zeta L(s,\chi_{-23})L(s,\rho)^2.
\]
To justify using the factorization on the right half-plane, ordinary
RH makes $\zeta$ nonzero there away from its pole, so
$L(s,\rho)=\zeta_K(s)/\zeta(s)$ is holomorphic there; the value at
one is obtained by removing the simple poles. The one-dimensional
factor is nonzero there by Dirichlet GRH. Consequently the remaining
zeros are exactly the degree-two target zeros, with their indicated
multiplicities. The functional equations finish both RH statements.

For comparison with a full prime-power computation, the squarefree
polynomial discriminant implies $\mathcal O_K=\mathbb Z[\alpha]$.
Modulo 23,
\[
 T^3-T-1=(T-3)(T-10)^2.
\]
There are two prime ideals above 23, each of residue degree one
(and ramification indices one and two), so the local zeta factor is
$(1-23^{-s})^{-2}$ and its logarithmic-derivative coefficient at
$23^j$ is $2\log23$ for every $j\ge1$. The prime-only statistic
omits 23 and requires no such local correction.

The root-count version also works for arbitrary irreducible
polynomials without assuming holomorphy of every irreducible Artin
factor: the Dedekind zeta function itself supplies the required
meromorphic continuation and nonnegative zero multiplicities.

\paragraph{What remains to be proved.}
The new task is an arithmetic upper bound for
\eqref{eq:sampling-polynomial-statistic}, uniformly over all sufficiently
large $n$. No finite computation establishes that bound. Positivity
and interpolation reduce the number of scales at which a proposed
arithmetic argument must work; they do not generate the needed
arithmetic information. To transfer the repository's determinant
construction, one must represent these root-count or splitting
weights by new arithmetic objects and prove their shortness,
spanning, and determinant lower bounds with the required error.
